% ============================================================
% Triple-Barrier Method Figure
% CSC491 Textbook, Chapter: Machine Learning
%
% Preamble requirements:
%     \usepackage{tikz}
%     \usepackage{pgfplots}
%     \pgfplotsset{compat=1.18}
%     \usetikzlibrary{patterns, fillbetween}
%     \usepackage{xcolor}
%     \definecolor{kuroseblue}{RGB}{30,  90, 160}
%     \definecolor{barrierred}{RGB}{180, 40,  40}
%     \definecolor{barriergray}{RGB}{100,100,100}
% ============================================================

\begin{figure}[ht]
\centering
\begin{tikzpicture}
\begin{axis}[
  width=0.82\textwidth,
  height=7cm,
  axis lines=left,
  xlabel={Time (bars)},
  ylabel={Price},
  xmin=0, xmax=10.5,
  ymin=93,  ymax=118,
  xtick=\empty,
  ytick={97, 100, 113},
  yticklabels={Stop-loss, Entry, Profit target},
  tick label style={font=\small},
  label style={font=\small},
  clip=false,
]

  %% --- Shaded corridor between barriers ---
  \fill[black!5] (axis cs:0,97) rectangle (axis cs:8,113);

  %% --- Upper barrier (profit) ---
  \addplot[
    dashed,
    color=barrierred,
    line width=1.2pt,
    domain=0:8,
  ] {113};
  \node[anchor=west, font=\small, text=barrierred]
    at (axis cs:8.1, 113) {Profit barrier (+1)};

  %% --- Lower barrier (stop-loss) ---
  \addplot[
    dashed,
    color=barriergray,
    line width=1.2pt,
    domain=0:8,
  ] {97};
  \node[anchor=west, font=\small, text=barriergray]
    at (axis cs:8.1, 97) {Stop-loss barrier (0)};

  %% --- Time barrier (vertical) ---
  \addplot[
    dashed,
    color=black!45,
    line width=1.0pt,
  ] coordinates {(8,93) (8,118)};
  \node[anchor=south, font=\small, text=black!55, rotate=90]
    at (axis cs:8, 107) {Time expiration};

  %% --- Price path: dips toward stop-loss, recovers, hits profit barrier ---
  %% Realistic wiggly path to show the path MATTERS for labeling
  \addplot[
    smooth,
    color=kuroseblue,
    line width=1.6pt,
  ] coordinates {
    (0,100)
    (0.5,101.5)
    (1.0,100.8)
    (1.5,99.5)
    (2.0,98.2)    %% dips close to stop-loss — dangerous moment
    (2.5,98.8)
    (3.0,100.5)
    (3.5,102.0)
    (4.0,101.2)
    (4.5,103.5)
    (5.0,105.8)
    (5.5,107.0)
    (6.0,108.5)
    (6.5,110.2)
    (7.0,111.5)
    (7.4,113)     %% hits profit barrier
  };

  %% --- Entry point ---
  \addplot[
    only marks,
    mark=*,
    mark size=3pt,
    color=kuroseblue,
  ] coordinates {(0,100)};
  \node[anchor=north east, font=\small, text=kuroseblue]
    at (axis cs:0,100) {Entry};

  %% --- Near-miss annotation (dip toward stop-loss) ---
  \draw[->, black!50, line width=0.8pt]
    (axis cs:2.8,96.5) -- (axis cs:2.1,98.0);
  \node[anchor=north, font=\scriptsize, text=black!55, align=center]
    at (axis cs:2.8,96.5) {Near stop-loss\\(fixed-horizon misses this)};

  %% --- Hit point on profit barrier ---
  \addplot[
    only marks,
    mark=*,
    mark size=3pt,
    color=barrierred,
  ] coordinates {(7.4,113)};
  \node[anchor=south west, font=\small, text=barrierred]
    at (axis cs:7.4,113) {\textbf{Label = 1 (profitable)}};

\end{axis}
\end{tikzpicture}

\caption{%
  The Triple-Barrier Method. A trade opens at the entry point (blue dot) and the
  price path is tracked bar-by-bar within a corridor bounded by a profit barrier
  (red, dashed), a stop-loss barrier (gray, dashed), and a time expiration
  (vertical, dashed). The label $y$ is assigned based on whichever barrier the
  price hits first. Here, the price dips dangerously close to the stop-loss before
  recovering and eventually hitting the profit barrier---earning label $y = 1$.
  A fixed-horizon method would assign the same label without ever ``seeing'' the
  near-miss, illustrating why the path matters for realistic trade labeling.%
}
\label{fig:triple-barrier}
\end{figure}